\documentclass[pdflatex,sn-vancouver-num]{sn-jnl}

%
% Set to \finalsubmissiontrue when compiling on a system without Inkscape and/or
% -shell-escape (e.g. most cloud/journal LaTeX services). Figures are then taken
% from pre-rendered PDFs instead of being converted from SVG on the fly; run
% convert-svgs-to-pdf.sh once beforehand to produce them.
\newif\iffinalsubmission
\finalsubmissionfalse

\usepackage{graphicx}%
\iffinalsubmission\else\usepackage{svg}%
\fi
\usepackage{multirow}%
\usepackage{amsmath,amssymb,amsfonts}%
\usepackage{amsthm}%
\usepackage{mathrsfs}%
\usepackage[title]{appendix}%
\usepackage{xcolor}%
\usepackage{textcomp}%
\usepackage{booktabs}%
\usepackage{algorithm}%
\usepackage{algorithmicx}%
\usepackage{algpseudocode}%
\usepackage{listings}%
\usepackage{csquotes}
\usepackage{siunitx}
\usepackage{todonotes}
\usepackage{array}
\usepackage{csvsimple}
\usepackage{etoolbox}
\usepackage{url}
\DeclareMathOperator*{\argmax}{arg\,max}

%
% Figures float more freely than the class allows by default; the same settings the
% other papers of this folder use.
\renewcommand{\topfraction}{0.9}
\renewcommand{\bottomfraction}{0.8}
\renewcommand{\textfraction}{0.07}
\renewcommand{\floatpagefraction}{0.75}

%
% A value that has not been measured yet renders as a visible marker rather than as
% a blank or a plausible-looking number. Everything numeric in this paper is meant to
% come from a generated file, as it does in ft-paper: see \dbstat below.
\providecommand{\sysUnknown}{\textbf{??}}

%
% Looks up a generated statistic. Values live in results/*.tex, written by whatever
% script ends up producing them, and are never typed into the prose. An unknown key
% renders as \sysUnknown so that a missing measurement is visible in the PDF instead
% of silently absent.
%   \dbstat{<kind>}{<key>}
\newcommand{\dbstat}[2]{\ifcsname dbstat@#1@#2\endcsname\csname dbstat@#1@#2\endcsname\else\sysUnknown\fi}

% Pulled in only when it exists, so that the paper builds before any measurement does.
\IfFileExists{results/smoothystats.tex}{\input{results/smoothystats.tex}}{}

\begin{document}

    \title{Smoothy: Using Topic Modeling for abundance estimation from metagenomic results}

    \author*[1]{\sur{Daniel} \fnm{Pfeifer}}\email{daniel.pfeifer@hs-heilbronn.de}

    \affil*[1]{\orgdiv{IT Faculty}, \orgname{Heilbronn University}, \orgaddress{\street{Max-Planck-Str. 39}, \city{Heilbronn}, \postcode{74081}, \state{Baden-W\"urttemberg}, \country{Germany}}}

    \date{September 2026}

    \abstract{\textbf{Background:}

    \textbf{Results:}

    \textbf{Conclusions:} }

    \keywords{Metagenomics, abundance estimation, topic modeling, latent Dirichlet allocation, read classification, taxonomic profiling}

    \maketitle

    \section{Introduction}\label{introduction}
    \todo{What abundance estimation has to do that classification does not, and why the assignments a
    classifier produces are not already an answer. The gap to state precisely: a read that a database
    leaves at an internal node names no organism, and the share of such reads is not small --- in the
    clinical runs of the companion paper a tenth of the classified reads sit at a genus, and the node
    they sit at leaves anywhere from eight to a hundred and fifty-seven species open.}

    \section{Related work}\label{relatedwork}
    Probabilistic topic modeling begins with probabilistic latent semantic analysis
    \citep{hofmann1999plsa}, which fits a document as a mixture over topics by maximum likelihood,
    and with latent Dirichlet allocation \citep{blei2003lda}, which places the mixtures under a
    Dirichlet prior and makes the model generative. Maximum-a-posteriori estimation of that model was
    derived by \citet{defreitas2001bayesian} and \citet{chien2008adaptive}; it is Hofmann's fit with
    the prior entering as a pseudo-count. Of the inference methods that followed, collapsed Gibbs sampling
    \citep{griffiths2004scientific} is the one this paper reaches for; \citet{asuncion2009smoothing}
    compare it against the others systematically and give the hyperparameter re-estimation step built
    on \citet{minka2000dirichlet} that is known as Minka's update.
    \citet{pfeifer2019topicgrouper} set out the same correspondence between documents, words and
    topics that this paper carries over to samples, taxonomy nodes and genomes.

    \todo{The other half of the related work, which is abundance estimation itself. The paper has to
    say what it is not: the re-estimation approaches that redistribute ambiguous assignments down the
    taxonomy, the expectation-maximisation family that fits a mixture of reference genomes to a read
    set, and the marker-gene profilers that sidestep the problem by only counting reads they can
    place. What topic modeling offers that none of them do is the thing to name here.}

    \section{Method}\label{method}
    \todo{This paper leads the implementation rather than reporting it, so this section is the
    specification: complete enough to build from, and precise enough that a measurement can
    contradict it.}

    \subsection{The generative model}\label{model}
    Shotgun sequencing draws reads from the whole of a genome rather than from one locus, so to a
    first approximation every $k$-mer of an organism's genome is equally likely to turn up in the
    reads. That is the assumption the method rests on, and it is why the method does not apply to
    metabarcoding data.

    Under it, the per-taxon counts a classifier reports are a mixture. Let $G$ be the sequenced
    genomes the database was built from, $N$ the nodes of the taxonomy it uses, and $c_n$ the
    $k$-mer count that the analysis of a sample $s$ reports at node $n$. Let $p(n \mid g)$ be the
    probability that a $k$-mer drawn from genome $g$ is stored at node $n$ of that database, and
    $p(g \mid s)$ the share of $g$ in the sample --- the abundance we are after. A $k$-mer observed
    at $n$ may have come from any genome, so
    %
    \begin{equation}\label{smoothylikelihood}
        \widehat{p}(g \mid s) \;=\; \argmax_{p(g \mid s)} \;\prod_{n \in N}
            \Bigl( \sum_{g \in G} p(n \mid g)\, p(g \mid s) \Bigr)^{c_n}.
    \end{equation}
    %
    The mixing is worst high in the tree, where a single node is fed by many genomes at once. That is
    exactly where a classifier's own answer says least, which is what the model has to offer: it
    reads the counts at the ambiguous nodes as evidence rather than discarding them.

    A match result reports two counts per taxon --- the $k$-mers attributed to it and the reads
    classified to it --- and $c_n$ is the former. Both the reason for that and the cost of it are
    worth stating.

    The cost is that the $k$-mers of one read overlap and are therefore far from independent: a read
    of length $L$ contributes $L - k + 1$ of them, and the likelihood counts what is really one
    observation as hundreds. This does not move the maximiser, since a factor common to every $c_n$
    cancels out of Equation \eqref{smoothylikelihood}, but it overstates how much evidence there is
    wherever the magnitude of the counts is read rather than their ratios --- against a prior, in any
    interval derived from the posterior, and most of all in a sample with few reads to begin with,
    where a single read landing on an unusual locus can carry a node by itself.

    Against that, $k$-mer counts use the whole of a read rather than a summary of it. A read informs
    several nodes at once --- its species-specific $k$-mers one, those it shares with a neighbour
    another, its core $k$-mers a third --- and the ambiguous part, which is the part this method
    exists to use, is kept rather than resolved away.

    Read counts are not the alternative they appear to be, and the objection is not that they are
    coarser. A read is placed at one node by the classifier's own rule, and that rule is built to name
    the most specific taxon the read supports, so it assigns systematically deeper than the nodes at
    which the read's $k$-mers actually sit --- for good reasons, since naming a genus where a species
    is warranted is the failure a classifier is designed to avoid. But $p(n \mid g)$ says where the
    database \emph{stores} the $k$-mers of $g$. It is not a statement about where a classifier would
    place a read of $g$. Pairing read counts with it would put an observation model and an emission
    distribution together that describe different processes, and the discrepancy is a systematic
    skew towards the deep nodes rather than noise, so it would not average away with more data.

    Using read counts would therefore require a different $p(n \mid g)$ altogether: the probability
    that a \emph{read} of $g$ is classified at $n$, which depends on the classifier's rule, on read
    length and on error rate. That is a property of the analysis and the sequencing run rather than
    of the database, so it could no longer be computed once and reused --- and it would have to be
    obtained by simulation, since nothing in the database implies it.

    The two inputs come from different places, and the difference matters. The $c_n$ describe one
    sample and are what a tool such as Genestrip already reports. The $p(n \mid g)$ describe the
    \emph{database} and not the sample: they are fixed once the database is built, so they are
    computed once and reused for every sample analysed against it.

    This is the correspondence with document-oriented topic modeling
    \citep{hofmann1999plsa,blei2003lda}, and it is close. A sample $s$ plays the part of a document,
    a taxonomy node $n$ that of a word, and a sequenced genome $g$ that of a topic; $c_n$ is then the
    count of a word in a document and $p(n \mid g)$ the topic--word distribution, in the roles set out
    in \citet{pfeifer2019topicgrouper}. Where the analogy stops is the part that makes the problem easy. Fitting a topic
    model to a corpus means learning the topics and the per-document mixtures together, and it is the
    first of those that makes the objective non-convex and the answer dependent on where the fit was
    started. Here the topics are not learned at all: $p(n \mid g)$ follows from the database, which is
    known before any sample is seen. What remains is only the per-document mixture $p(g \mid s)$ ---
    the analogue of inference for an unseen document against an already-fitted model, not of fitting
    one.

    That is what makes Equation \eqref{smoothylikelihood} well behaved. With the $p(n \mid g)$ fixed,
    its logarithm is a positive combination of logarithms of functions linear in $p(g \mid s)$, hence
    concave over the simplex. There are therefore no local optima to be trapped in, and an iteration
    reaches the maximum wherever it starts. Whether that maximum is a single point is a separate
    question, and it turns on the taxa being distinguishable from one another; Section \ref{pgn}
    takes it up, since it is a property of the $p(n \mid g)$ rather than of the objective.

    \todo{Confirm the notation. The idea as given had $p(g \mid n)$ in the first factor; written that
    way the bracket is an inner product of two distributions over $G$ and is not a probability of
    anything, and it does not sum to one over $N$. With $p(n \mid g)$ the bracket is $p(n \mid s)$,
    the likelihood is a proper multinomial over the observed counts, and the quantity the new goal
    measures --- the $k$-mers of $g$ that the database stores at $n$, normalised over $n$ --- is
    exactly $p(n \mid g)$. Equation \eqref{smoothylikelihood} is written accordingly.}

    \todo{Confirm the correspondence as written --- sample to document, node to word, genome to topic
    --- and that only $p(g \mid s)$ is inferred, the topics being given by the database rather than
    learned from the samples.}

    \subsection{Estimating \texorpdfstring{$p(n \mid g)$}{p(n|g)}}\label{pgn}
    The topic--word distribution is a property of the database, so it is measured on the database and
    not fitted to anything. A new Genestrip goal walks the genomes the update phase considers ---
    which is every genome of the release in scope, and not only those the entry phase stored --- and
    for each of them looks up every $k$-mer in the database $d$.

    Two cases arise, and the second is the one that matters. If the genome $g$ was entered into $d$,
    a $k$-mer of $g$ that $d$ holds at node $n$ raises the counter $c(g,n)$. If $g$ was not entered,
    there is no row to raise: a $k$-mer of $g$ that $d$ holds at $n$ raises that node's
    \textsc{other} counter $c(\textsc{other},n)$ instead, and a $k$-mer that $d$ does not hold at all
    raises nothing. Counting is by frequency rather than by distinct $k$-mer, since a $k$-mer shared
    between genomes is the evidence the mixture is about and discarding its multiplicity would
    discard the ambiguity the model exists to resolve.

    A row is kept for every taxon and not only for the taxa that carry a genome. For a data taxon the
    row is that genome's. For a taxon $t$ above the data, every genome subordinate to $t$ contributes
    its $k$-mers to $t$ and to $t$'s ancestors, and only there: a $k$-mer of such a genome counts
    towards $t$'s row when the database stores it at $t$ or at a node above $t$, and not when it is
    stored below. So $c(t,n) = \sum_{g \preceq t} c(g,n)$ for $n$ on the path from $t$ to the root,
    which is the same restriction the storage layout below rests on: a row never reaches beneath its
    own taxon.

    So at a rank above the data taxa a topic is not an organism. Inferring at genus rank returns a
    distribution over genera, and the row read for a genus is the one written at its own node: the
    pooled contribution of every genome beneath it. The topic is that clade's pan-genome as the
    reference collection happens to represent it --- and, counting being by frequency, a pan-genome
    weighted towards whatever its members were sequenced most often. A $k$-mer carried by many
    genomes of the clade weighs more than one carried by a single draft.

    That weighting is defensible, since a $k$-mer common to a clade is likelier to be observed from
    it, but it makes the topic a property of the reference collection as much as of the organisms.
    Where a clade is dominated by one heavily sequenced member --- \emph{S.~pneumoniae} accounts for
    more than nine thousand of the twenty-two thousand \emph{Streptococcus} assemblies in RefSeq ---
    the genus topic is largely that member's genome wearing the genus's name.

    The estimate is then
    %
    \begin{equation}\label{smoothypng}
        p(n \mid g) \;=\; \frac{c(g,n)}{\lvert K(g) \rvert},
    \end{equation}
    %
    where $K(g)$ is the set of $k$-mers of $g$ --- all of them, not only those $d$ holds.

    That denominator is deliberate and has a consequence worth being explicit about: the row
    $p(\cdot \mid g)$ does not sum to one. What is missing is the share of $g$'s $k$-mers the database
    does not hold anywhere, and leaving it missing is what makes a genome that $d$ covers thinly
    weigh less than one it covers well. The observed counts $c_n$ are subject to the same omission,
    since a $k$-mer absent from $d$ can never be reported at a node either.

    The \textsc{other} row plays the part of an organism the database does not contain. Without it
    every count would have to be explained by some $g \in G$, and a sample holding an organism absent
    from $d$ would have its abundance pushed onto whichever genomes happen to resemble it. It is the
    same role the \textsc{other} bucket plays in a refined tree. Its $k$-mers come from many genomes
    at once, so its denominator pools them:
    %
    \begin{equation}\label{smoothyother}
        p(n \mid \textsc{other}) \;=\;
            \frac{c(\textsc{other},n)}{\sum_{g \notin d} \lvert K(g) \rvert},
    \end{equation}
    %
    the sum running over every genome the entry phase did not store. \textsc{other} is therefore a
    genome-weighted average over everything the database left out, rather than a quantity belonging
    to any one organism.

    Every taxon carries a row, the ones above the data taxa included, but a single inference does not
    use them all. A rank is chosen and held fixed for the run, and $G$ in
    Equation \eqref{smoothylikelihood} is then the set of taxa at that rank. The same stored
    fractions serve every rank, since a row is written for each taxon regardless: choosing the rank
    is choosing which rows to read, not recomputing them.

    \begin{table}[!htbp]
        \centering
        \footnotesize
        \begin{tabular}{@{}l|l|l@{}}
            \textbf{Topic modeling} & \textbf{Here} & \textbf{Symbol}\\
            \hline
            document $j$ & sample (statistic) $s$ & $s$, or $s_1 \dots s_m$\\
            word $w$ & taxonomy node & $n \in N$\\
            topic $k$ & taxon at the chosen rank & $g \in G$\\
            word count in a document & count reported at a node & $c_n$, corrected $\hat{c}_n$\\
            topic--word distribution $\varphi_{wk}$ & where a taxon's $k$-mers are stored &
                $\varphi_{ng} = p(n \mid g)$\\
            document--topic distribution $\theta_{kj}$ & abundance in the sample &
                $\theta_g = p(g \mid s)$\\
            Dirichlet concentration & unchanged & $\alpha$\\
        \end{tabular}
        \caption{The correspondence of Section \ref{model} in the symbols each field uses. The
        indices differ deliberately: $k$ is the $k$-mer length throughout this paper and cannot also
        be a topic index, and $w$ for a node would say nothing. The distributions keep their usual
        names, $\varphi$ and $\theta$, so that the equations below read as they do in the literature
        the method is borrowed from.}\label{correspondence}
    \end{table}

    Fixing the rank is also what keeps the estimate identifiable. Were an ancestor and its
    descendants topics at once, the same counts could be explained by either and no data could
    separate them; at one rank that competition does not arise, because no topic is an ancestor of
    another.

    One degeneracy survives it, and it is not hypothetical. Two taxa of the chosen rank are
    distinguishable only where their columns differ, and a taxon that retains none of its own
    $k$-mers has a column supported entirely on the ancestors it shares with its siblings. Where two
    such taxa sit under the same parent, their columns are proportional and any split of abundance
    between them fits the data equally well. This is the mitis complex of the companion paper seen
    from the estimation side: in a \emph{Streptococcus} database built on the reference taxonomy
    \emph{S.~mitis} retains nothing at all, so at species rank it cannot be told from the
    pneumococcus by any amount of sequencing. The condition is a property of the database and can be
    detected from the stored fractions before a sample is seen, which is the right place to report it
    rather than leaving a reader to wonder why two abundances trade off.

    \todo{Check the degeneracy argument above before relying on it: that two taxa of the chosen rank
    are indistinguishable exactly when their columns are proportional, and that a taxon retaining
    none of its own $k$-mers is the way this happens in practice. If it holds, the implementation can
    report the affected taxa from the stored fractions alone, before any sample is analysed.}

    \todo{Whether the pooling should be weighted by sequencing effort at all, or whether a clade's
    genomes should enter its row with equal weight regardless of how many times the clade was
    sequenced. As written the first holds, which is what the counting gives for free; the second
    would make the topic a property of the clade rather than of the collection, at the cost of an
    assumption about which genomes represent it. Decide it, and say which, because the two give
    different abundances for exactly the clades a reader cares most about.}

    \todo{Which rank to run at is a decision the paper has to give guidance on, not only a parameter.
    A deep rank buys resolution the data may not support; a shallow one throws away evidence the
    database does hold. The companion paper's measurements of how far down the tree a sample's
    evidence actually reaches are the natural place to take that guidance from.}

    \todo{Whether Equation \eqref{smoothypng} should be renormalised over $n$: as written the
    likelihood of Equation \eqref{smoothylikelihood} is not a proper multinomial, which is
    defensible --- the missing mass is unobservable and the weighting it induces is the one we want
    --- but it should be argued rather than passed over.}

    \todo{A consequence of the pooled denominator to check on data rather than to assume away. A
    stored genome's row nearly sums to one, since the entry phase put its $k$-mers into $d$; the
    \textsc{other} row does not come close, because most $k$-mers of most unstored genomes are absent
    from a database built around one clade. In the mixture, a row of small total mass has to be
    carried by a correspondingly large $p(g \mid s)$ to explain the same counts. So a sample holding
    an organism the database lacks should still be recognised as such --- nothing else can explain
    its counts --- but the size of $p(\textsc{other} \mid s)$ will not be on the same scale as the
    abundances beside it, and reporting it as though it were would misstate how much of the sample is
    unaccounted for. This belongs to Section \ref{abundances}.}

    \paragraph{Cost and storage.} The goal runs once for a database and writes its result to a
    serialized file, which every later analysis against that database reads. Nothing here is paid per
    sample.

    The result is small, and for a reason that belongs to the method rather than to the encoding. A
    $k$-mer of $g$ is stored by the LCA update at the lowest common ancestor of the taxa that carry
    it; $g$ is one of them, so that node is an ancestor of $g$'s own taxon or the taxon itself. Every
    other node is therefore unreachable for $g$, and $p(n \mid g)$ is zero there by construction. The
    row for $g$ needs no entry for it. What is stored is one array of doubles per taxon, holding the
    fractions along the path from that taxon up to the root --- an entry for the taxon itself and one
    for each ancestor, so a path of length $n$ costs $n$ doubles.

    That turns a matrix over genomes and nodes into one over genomes and tree depth. A bacterial
    taxonomy is some ten to fifteen deep where the data taxa sit, so a database entering a few
    thousand genomes stores a few hundred kilobytes rather than the tens of megabytes a dense matrix
    over all its nodes would need.

    \todo{The \textsc{other} row is the exception and needs its own encoding. Its $k$-mers come from
    genomes the database never stored, so the node a $k$-mer of theirs is found at is the ancestor of
    whichever \emph{stored} taxa carry it and bears no relation to any path of its own.
    \textsc{other} therefore ranges over the whole taxonomy and does not compress the way the rows
    above it do --- one entry per node that carries anything, kept sparsely.}

    \todo{What the pass itself costs. One traversal per genome of the update scope against a built
    database. For a bacterial project that scope is the whole division, so this is the same order of
    work as the LCA update rather than a small addition to it. State it honestly and measure it
    against the update, since it is the one part of the method that is not cheap.}

    \subsection{Correcting the counts for sequencing error}\label{error}
    A $k$-mer is observed only if all $k$ of its bases were read correctly. With a per-base error rate
    $e_b$ and errors taken as independent, an instance survives with probability $(1 - e_b)^k$, so the
    reported counts are a systematic undercount and the corrected count is
    %
    \begin{equation}\label{smoothyerror}
        \hat{c}_n \;=\; \frac{c_n}{(1 - e_b)^k}.
    \end{equation}
    %
    The factor is not a small adjustment. At $k = 31$ it is $0.73$ for $e_b = 1\,\%$, $0.53$ at
    $2\,\%$ and $0.20$ at the $5\,\%$ of a nanopore run: four $k$-mers in five are lost before
    anything is counted.

    Two things follow, and the first is the one to say plainly. Applied as written,
    Equation \eqref{smoothyerror} does not change the estimate at all. It scales every $c_n$ by the
    same factor, the objective $\sum_n c_n \log (\varphi\theta)_n$ scales with it, and the maximiser over
    $\theta$ is unchanged. The correction earns its place elsewhere: wherever the magnitude of the
    counts is read rather than only their ratios --- against a prior, whose weight is otherwise
    silently multiplied by the error rate; in any absolute quantity reported alongside the
    distribution; and in comparing samples sequenced at different error rates, where uncorrected
    counts are not on one scale.

    The second is that $e_b$ need not be a constant. The quality strings carry a per-base estimate,
    so the surviving fraction can be averaged over the reads actually seen instead of assumed. That
    matters most where it is largest: nanopore error varies widely between reads, and a mean rate
    understates the loss because $(1 - e_b)^k$ is convex in $e_b$.

    \paragraph{Counts that the assumption does not fit.} The correction above rescales every node
    alike and therefore says nothing about whether the assumption behind
    Equation \eqref{smoothylikelihood} holds at any particular one. There is a measure that does, and
    it comes from the classifier itself. The $k$-mer consistency ratio of \citet{genestrip-pub}
    compares the distinct $k$-mers actually seen for a taxon against the number expected if its
    $k$-mers were drawn uniformly: with $u_{t,d}$ the $k$-mers the database stores for $t$,
    $k_{t,f,d}$ the count observed and $u_{t,f,d}$ the distinct ones among them,
    %
    \begin{equation}\label{smoothyconsistency}
        E \;=\; u_{t,d} \Bigl( 1 - (1 - 1/u_{t,d})^{k_{t,f,d}} \Bigr),
        \qquad r_{t,f,d} \;=\; u_{t,f,d} \,/\, E .
    \end{equation}

    Its assumption is worth reading twice, because it is ours. It holds where every $k$-mer of a taxon
    is equally likely to be seen and fails on repeats, on genomic dust and under amplification
    protocols --- which is word for word the assumption Section \ref{model} rests on and the reason
    this method does not apply to metabarcoding. So $r$ is not an external criterion bolted on. It is
    the method's own premise, tested against the data, per node and per sample.

    That makes it a weight rather than a filter. A threshold on $r$ would discard a node's counts
    whole and turn a continuous doubt into a binary one; the same information is better spent
    attenuating a node's contribution in proportion to how far its counts depart from what the
    assumption predicts. It enters where the error correction does,
    %
    \begin{equation}\label{smoothyweight}
        \hat{c}_n \;=\; \frac{w(r_n)\, c_n}{(1 - e_b)^k} ,
    \end{equation}
    %
    and unlike that correction it is not a common factor: $w$ varies from node to node, so this one
    does move the estimate rather than only its scale. That is the point of it.

    A form that meets the requirements below, and the one we propose, is
    %
    \begin{equation}\label{smoothyw}
        w(r) \;=\; \max\Bigl( 0, \; 1 + \lambda \ln \bigl( \min(r, 1) \bigr) \Bigr),
        \qquad \lambda > 0 .
    \end{equation}
    %
    The inner $\min$ discards the uninformative direction, the logarithm makes the penalty gentle near
    one and steep as $r$ falls, and the outer $\max$ puts a floor at zero: a node below
    $r = e^{-1/\lambda}$ contributes nothing. Filtering is therefore not rejected but recovered, as
    the limiting case of a weight rather than as a separate rule with a threshold of its own.

    One parameter remains, and $\lambda$ says how severely inconsistency is punished. It is worth
    noting that a calibration point already exists: the worked example of \citet{genestrip-pub} reports
    $r = 0.8631$ for Epstein--Barr virus in a saliva sample and reads it as \emph{supporting} the
    consistency of the counts. A node so described should keep most of its weight, which
    $\lambda = 1$ does not do --- it leaves $0.85$ --- while $\lambda = 1/2$ leaves $0.93$. That
    suggests $\lambda$ well below one, but it is one example and the value wants measuring.

    What $w$ must satisfy, however $\lambda$ is fixed:
    %
    \begin{itemize}
        \item $w(1) = 1$. A node whose counts match the expectation is not penalised.
        \item $w$ falls as $r$ falls below one, and smoothly. A step function is a filter wearing a
              weight's clothes, with a tuning parameter and a cliff.
        \item $w \approx 1$ for $r$ above one. Seeing appreciably more distinct $k$-mers than expected
              is not a thing that happens, so that direction carries no information and should not be
              read as one.
        \item $w$ depends on the count as well as the ratio. At small $k_{t,f,d}$ the ratio is noise,
              and a node holding three $k$-mers must not be discounted because its $r$ came out at
              $0.6$. Equation \eqref{smoothyw} does not yet do this, and the natural repair is to
              shrink $r$ towards one as the count falls --- replacing $r$ by
              $1 + (r - 1)\,k_{t,f,d}/(k_{t,f,d} + k_0)$ for some scale $k_0$ --- rather than to
              complicate $w$ itself.
        \item $w$ is a heuristic in $r$ and not a tail probability. The exact distribution of
              $u_{t,f,d}$ is known and approximately normal about $E$, but it is so sharply peaked
              that small deviations drive it to zero, which is why \citet{genestrip-pub} report the
              ratio rather than a test.
    \end{itemize}

    \todo{Fix $\lambda$ and $k_0$ against data rather than by argument, and check
    Equation \eqref{smoothyw} against alternatives while doing so. The measurements that decide it:
    how $r$ is distributed across nodes in real samples, and whether the nodes with low $r$ are the
    ones an independent standard says are wrong.}

    \todo{The effect that would change the estimate is the one not modelled here, and the section
    should say so rather than leave it implied. A corrupted $k$-mer does not only disappear: it can
    still be in the database, at some other node. How often depends on how densely the database
    populates the neighbourhood of the original $k$-mer, and that is not uniform --- a $k$-mer shared
    across a clade sits in a denser neighbourhood than one specific to a strain. Errors would then
    move counts between nodes rather than merely shrink them, which no scalar correction can undo.
    Measure whether it matters before modelling it: mutating the $k$-mers of a simulated read set at a
    known rate and re-classifying gives the size of the effect directly.}

    \subsection{Priors from the taxonomy}\label{priors}
    \todo{A metagenomic problem differs from a text one in having a reference taxonomy, so the model
    need not learn from the data what is already known. Where that knowledge enters, and what it
    costs when the taxonomy is wrong. This is not the symmetric Dirichlet of
    Appendix \ref{prior} --- that prior knows nothing about the tree and treats every taxon as
    exchangeable, which is precisely what a taxonomy says they are not. If anything here survives, it
    is because it uses structure the Dirichlet throws away.}

    \subsection{Inference}\label{inference}
    Write $\theta_g = p(g \mid s)$, $\varphi_{ng} = p(n \mid g)$ and $C = \sum_n \hat{c}_n$, following
    the usual symbols of the topic-modeling literature for the mixture and for the topic--word
    distribution. The estimate
    maximises $\ell(\theta) = \sum_n \hat{c}_n \log (\varphi\theta)_n$ over the simplex. Two ways to
    compute it are worth setting out, because the second follows from the first and needs less
    apparatus.

    \paragraph{Constrained ascent.} With the constraint in affine form the Lagrangian is
    %
    \begin{equation}\label{smoothylagrange}
        \Lambda(\theta, \lambda) \;=\; \sum_{n \in N} \hat{c}_n \log (\varphi\theta)_n
            \;-\; \lambda \Bigl( \sum_{g \in G} \theta_g - 1 \Bigr),
    \end{equation}
    %
    and stationarity gives $\sum_n \hat{c}_n \varphi_{ng} / (\varphi\theta)_n = \lambda$ at every topic carrying
    mass. The multiplier is not left to be solved for: multiplying that condition by $\theta_g$,
    summing over $g$ and using $\sum_g \varphi_{ng} \theta_g = (\varphi\theta)_n$ collapses the left side to
    $\sum_n \hat{c}_n$, so
    %
    \begin{equation}\label{smoothylambda}
        \lambda \;=\; C .
    \end{equation}
    %
    Equation \eqref{smoothylagrange} handles $\sum_g \theta_g = 1$ and not $\theta_g \geq 0$, which
    is the constraint that binds: most taxa are absent from any one sample, so the maximiser lies on
    the boundary of the simplex and the stationarity condition holds there as the inequality
    $\sum_n \hat{c}_n \varphi_{ng} / (\varphi\theta)_n \leq C$. Ascent therefore needs either the full
    Karush--Kuhn--Tucker conditions, or a projection of each iterate back onto the simplex, or a
    mirror step in the Kullback--Leibler geometry that keeps positivity by construction. All three
    work, and all three need a step size.

    \paragraph{Expectation--maximisation.} Substituting $\lambda = C$ into the stationarity condition
    and reading it as a fixed point gives the iteration directly, with no step size and no projection.
    It is the mixture-proportion step of \citet{hofmann1999plsa} with the topic--word distributions
    held fixed:
    %
    \begin{equation}\label{smoothyem}
        \theta_g \;\longleftarrow\; \frac{\theta_g}{C} \sum_{n \in N}
            \frac{\hat{c}_n \, \varphi_{ng}}{(\varphi\theta)_n}.
    \end{equation}
    %
    So the two are not rivals but two readings of one condition. The quantity
    $\varphi_{ng}\theta_g / (\varphi\theta)_n$ is the share of the count at $n$ that topic $g$ is responsible for;
    the step reassigns each $\hat{c}_n$ among the topics in that proportion and renormalises. Four
    properties make it the one to implement.

    It needs no constraint machinery. The update sums to one over $g$ by construction, because the
    responsibilities do, so the simplex is preserved exactly and no multiplier or projection appears.
    Non-negativity is preserved for the same reason. It has no step size, so nothing has to be tuned
    or line-searched. It increases the likelihood at every step, by the usual
    minorise--maximise argument, and that argument needs only the responsibilities to be a
    distribution over $g$ --- which holds whether or not the rows of $\varphi$ are normalised over $n$, so
    the question left open in Section \ref{pgn} does not affect the algorithm. And with the objective
    concave, the point it converges to is the global maximum rather than a stationary point that
    happens to be near where it started.

    Each iteration costs one pass over the non-zeros of $\varphi$. Those are few: a column is supported on
    the path from its taxon to the root, so $\varphi$ has of the order of one entry per taxon per level
    rather than one per taxon per node.

    \todo{Two known weaknesses to state rather than discover. EM converges linearly, and slowly when
    two topics are near-collinear --- which is exactly the degenerate case of Section \ref{pgn}, so
    slow convergence is a symptom worth reporting rather than only a cost. And the update is
    zero-preserving: a topic that starts at zero stays there, so the iteration is started uniform.
    If the sizes turn out to warrant it, SQUAREM-style acceleration or a few Newton steps on the
    support at the end are the usual remedies; measure before adding either.}

    \todo{Exact sparsity, if it is wanted. EM drives absent topics towards zero but reaches it only in
    the limit, so a report will list every taxon with a small positive abundance unless something is
    done. A Dirichlet prior with concentration below one turns the step into a MAP update of the same
    shape and gives genuine zeros; thresholding does not, it only hides them.}

    \subsection{From topics to abundances}\label{abundances}
    What Equation \eqref{smoothylikelihood} estimates is a share of $k$-mers, and therefore a share of
    \emph{sequence}. It is not a share of organisms. A genome twice as long contributes twice the
    $k$-mers at equal cell count, so reporting $\theta_g$ as an abundance overstates the large genomes
    and understates the small ones by exactly their length ratio. The two quantities are both worth
    reporting and must not be confused, which in practice they often are.

    Relative abundance follows by reweighting each topic by its genome length and renormalising:
    %
    \begin{equation}\label{smoothyabundance}
        a_g \;=\; \frac{\theta_g \,/\, \lvert K(g) \rvert}
                       {\sum_{g' \in G \setminus \{\textsc{other}\}}
                            \theta_{g'} \,/\, \lvert K(g') \rvert} ,
        \qquad g \in G \setminus \{\textsc{other}\} .
    \end{equation}
    %
    The length is measured in $k$-mers rather than bases, which is the same unit the counts are in and
    saves a conversion. It also costs nothing to obtain: $\lvert K(g) \rvert$ is already the
    denominator of Equation \eqref{smoothypng}, so the goal of Section \ref{pgn} has computed it
    before any sample is analysed and need only store it alongside the fractions.

    Above the data taxa the topic is a pan-genome and the length that belongs to it is the mean over
    the genomes pooled into it, $\overline{K}_t = \lvert G_t \rvert^{-1} \sum_{g \in G_t}
    \lvert K(g) \rvert$ with $G_t$ the genomes beneath $t$ --- the mean and not the sum that
    normalises its row, since what is wanted is the length of one organism of the clade and not the
    length of all of them together.

    \todo{Two limits of Equation \eqref{smoothyabundance} to state rather than let a reader find. The
    length is the reference genome's and not the sample organism's, so a strain carrying a large
    accessory complement is measured against the wrong ruler. And above the data taxa the mean is
    taken over an unevenly sampled set --- the same sequencing-effort bias that shapes the row itself,
    now entering the abundance a second time.}

    \paragraph{What the database cannot account for.} \textsc{other} is deliberately absent from
    Equation \eqref{smoothyabundance}. It is not an organism, it has no genome length worth the name,
    and dividing it by the mean length of everything the database omits would produce a number that
    invites comparison with the $a_g$ while meaning something else entirely. It is reported on its
    own instead, as
    %
    \begin{equation}\label{smoothyunaccounted}
        u \;=\; \theta_{\textsc{other}} ,
    \end{equation}
    %
    the share of the sample's \emph{sequence} that the database cannot name. The $a_g$ then sum to one
    over the taxa the database does know, and $u$ stands beside them rather than among them: a
    statement about the database's coverage of this sample, not about an organism in it.

    Its magnitude wants reading with care, and Section \ref{pgn} says why. Writing
    $\rho_g = \sum_n p(n \mid g)$ for the share of $g$'s $k$-mers the database holds anywhere, a
    stored genome has $\rho_g$ near one while \textsc{other} does not come close, since most $k$-mers
    of most omitted genomes are absent from a database built around one clade. A topic with small
    $\rho$ must carry a larger $\theta$ to explain the same counts, so $u$ overstates the unaccounted
    fraction by roughly $1/\rho_{\textsc{other}}$. Dividing each $\theta$ by its $\rho$ before
    comparing puts the topics back on one footing --- which is the same operation as normalising the
    rows of $P$, the question left open in Section \ref{pgn}. The two are one decision and should be
    taken together.

    \todo{Whether $u$ is reported raw or divided by $\rho_{\textsc{other}}$, and either way with the
    caveat stated. A reader will read it as `this fraction of the sample is something you do not have
    a genome for', and that reading has to be either true or explicitly corrected.}

    \subsection{What it costs}\label{cost}
    \todo{Time and space, stated before anything is measured, so that the measurement can be held
    against a prediction rather than merely reported.}

    \section{Results}\label{results}
    \todo{Nothing here until the method is implemented. Every number in this section must come from
    a generated file through \texttt{\textbackslash dbstat} or \texttt{\textbackslash csvreader};
    an unmeasured value shows as \sysUnknown\ rather than as a blank.}

    \todo{Abundance estimation has a ground truth that classification does not: a simulated read set
    knows the abundance it was drawn at. Say early which simulated sets are used and at which error
    rates, and keep at least one real sample where the answer is only known qualitatively, so that
    the method is not only measured where it is easiest to measure.}

    \section{Discussion}\label{discussion}
    \todo{Including what the method cannot do, which is the part a reader will look for.}

    \section{Conclusions}\label{conclusions}
    \todo{}

    \begin{appendices}

    \section{A Dirichlet prior and a sampler}\label{prior}
    Everything above estimates $\theta$ from the counts alone. A symmetric Dirichlet prior
    $\theta \sim \mathrm{Dir}(\alpha \mathbf{1})$, the scalar concentration of topic modeling, is the
    natural addition. It is set down here rather than in the method
    because it does not, on the evidence so far, earn its cost: the useful range of $\alpha$ is
    exactly the range in which maximising fails, so taking the prior means giving up the fixed point
    of Section \ref{inference} for a sampler, giving up concavity, and committing to analysing
    samples in sets rather than singly. That is a different method, not an option within this one.
    It is kept because the reasoning behind it is worth not repeating from scratch, and because the
    objection to it is contingent rather than absolute. That objection is the one of
    Section \ref{model}: $k$-mer counts overstate the evidence by roughly the read length, so a prior
    has nothing left to weigh against. Read counts would not rescue it --- Section \ref{model} gives
    the reason they are unusable at all --- but anything that restored an honest sample size would,
    whether that is thinning the $k$-mers of a read, an explicit overdispersion term, or a likelihood
    stated per read against a $p(n \mid g)$ built for one.

    \subsection{Where EM still works}\label{mapworks}
    Maximum-a-posteriori estimation of a topic model
    \citep{defreitas2001bayesian,chien2008adaptive} adds the prior to the maximum-likelihood fit of
    \citet{hofmann1999plsa} as a pseudo-count. For $\alpha \geq 1$ the objective is
    $\ell(\theta) + (\alpha - 1) \sum_g \log \theta_g$ and the iteration keeps its shape, gaining
    only a pseudo-count:
    %
    \begin{equation}\label{smoothymapem}
        \theta_g \;\longleftarrow\; \frac{\bigl( \sum_n \hat{c}_n \, \varphi_{ng} \theta_g /
            (\varphi\theta)_n \bigr) + \alpha - 1}{C + \lvert G \rvert (\alpha - 1)} .
    \end{equation}
    %
    The prior then shrinks the estimate towards the uniform distribution. That is smoothing, and it is
    the opposite of what a Dirichlet prior is usually wanted for here.

    \paragraph{Mode and mean differ by one.} The pseudo-count carries an offset that is easy to lose
    sight of when the two estimators of this appendix are compared. Equation \eqref{smoothymapem}
    converges to the posterior \emph{mode},
    $(m_g + \alpha - 1) / (M + \lvert G \rvert (\alpha - 1))$, while the sampler below reports the
    posterior \emph{mean}, $(m_g + \alpha) / (M + \lvert G \rvert \alpha)$. The two differ by exactly
    one in the concentration: the mode at $\alpha$ is the mean at $\alpha - 1$. Setting the same
    $\alpha$ for both and reading the difference in their output as a difference between maximising
    and sampling would therefore be a mistake --- part of it is only the offset, and at the small
    $\alpha$ of interest here that part is not small. It is also why the two have different admissible
    ranges: the mean is defined for every $\alpha > 0$ and the mode only from $\alpha = 1$ upwards.

    This agrees with the MAP update of \citet{asuncion2009smoothing}, whose estimate
    $\hat{\theta}_{kj} = (N_{kj} + \alpha - 1)/(N_j + K\alpha - K)$ is
    Equation \eqref{smoothymapem} in their notation, $K\alpha - K$ being $\lvert G \rvert (\alpha -
    1)$ here. They reach the same conclusion about the offset from the other direction, observing that
    MAP offsets by $-1$ against collapsed Gibbs sampling and that the two are brought into agreement
    by incrementing $\alpha$ by one --- which is why their own hyperparameter search shifts its grid
    by one for MAP, hyperparameters below one giving it negative probabilities. Only the mixture side
    of their derivation applies here: they estimate the topic--word distributions as well, and ours
    are given by the database.

    \subsection{Where it does not}\label{gibbs}
    Sparsity needs $\alpha < 1$, and there the MAP is not merely harder
    but ill-posed: the Dirichlet density is unbounded at the corners of the simplex, so the objective
    has no finite maximum and Equation \eqref{smoothymapem} produces negative numerators that have to
    be clamped, which is a repair rather than a solution. The remedy is to stop maximising over
    $\theta$ and integrate it out. Collapsed Gibbs sampling \citep{griffiths2004scientific} does that, and it is markedly cheaper here
    than in the corpus setting it comes from: the topics are fixed, so only the assignment of each
    counted $k$-mer occurrence is sampled and never a topic--word distribution. For an occurrence
    observed at node $n$,
    %
    \begin{equation}\label{smoothygibbs}
        p(z = g \mid \cdot) \;\propto\; \varphi_{ng} \, \bigl( m_g^{\neg} + \alpha \bigr),
    \end{equation}
    %
    with $m_g^{\neg}$ the number of occurrences currently assigned to $g$, that one excluded. A sweep
    visits every occurrence; the estimate is the posterior mean
    $\mathbb{E}[\theta_g] = (m_g + \alpha) / (M + \lvert G \rvert \alpha)$ over the retained sweeps,
    with $M$ the total number of occurrences.

    The path structure of Section \ref{pgn} bounds the work. An occurrence observed at node $n$ can
    only belong to a topic whose column has support at $n$, and $\varphi_{ng} > 0$ only where $n$ lies on
    the path from $g$ to the root --- that is, only for the rank-taxa lying beneath $n$. The candidate
    set is therefore small for a specific node and large only for the nodes near the root, which is
    the right shape: the sampler spends its time on exactly the occurrences that are ambiguous.
    Occurrences at one node are moreover exchangeable, so the state is the assignment counts per node
    rather than a label per $k$-mer, and a sweep re-draws a node's occurrences from the urn in turn
    instead of touching millions of tokens individually.

    \subsection{Estimating \texorpdfstring{$\alpha$}{alpha} needs more than one sample}
    Everything so
    far has concerned a single statistic $s$. A concentration parameter cannot be estimated from it:
    $\alpha$ describes how $\theta$ varies from one sample to the next, and one sample exhibits no
    variation. Fitting it anyway returns whatever value best explains the single $\theta$ at hand,
    which is not an estimate of anything.

    So the estimation is stated over a set. Let $s_1, \dots, s_m$ be statistics produced against the
    same database, $c^{(i)}_n$ the count that $s_i$ reports at node $n$, and
    $\theta^{(i)} = p(g \mid s_i)$ its mixture. The samples share the prior and nothing else:
    %
    \begin{equation}\label{smoothyjoint}
        p\bigl(c^{(1)}, \dots, c^{(m)} \bigm| \alpha\bigr) \;=\; \prod_{i=1}^{m} \int
            \mathrm{Dir}\bigl(\theta^{(i)} \bigm| \alpha \mathbf{1}\bigr)
            \prod_{n \in N} \bigl(\varphi\theta^{(i)}\bigr)_n^{\,c^{(i)}_n} \, d\theta^{(i)} .
    \end{equation}
    %
    Each $\theta^{(i)}$ stays a quantity of its own sample --- the aim is still one abundance
    distribution per sample, not a shared one --- while $\alpha$ is common to all $m$ and is what the
    set is pooled to determine. That is empirical Bayes, and it is why the method analyses samples in
    sets. A clinical cohort measured on one database is exactly such a set.

    Minka's update is the fixed point for $\alpha$: \citet{asuncion2009smoothing} apply the Dirichlet
    estimator of \citet{minka2000dirichlet} as a step of the fitting procedure and give it that name.
    It belongs to the sampler and not to the maximisation of Section \ref{mapworks}: the same authors
    leave MAP out of their hyperparameter-learning experiments precisely because Minka's update does
    not apply to it. So the two halves of this appendix do not combine into one procedure. Taking the
    prior means taking the sampler, and estimating $\alpha$ rather than fixing it means the same
    again. With $m_{ig}$ the
    occurrences of $s_i$ currently assigned to topic $g$ and $M_i = \sum_g m_{ig}$,
    %
    \begin{equation}\label{smoothyalpha}
        \alpha \;\longleftarrow\; \alpha \,
            \frac{\sum_{i=1}^{m} \sum_{g \in G} \bigl( \psi(m_{ig} + \alpha) - \psi(\alpha) \bigr)}
                 {\lvert G \rvert \sum_{i=1}^{m} \bigl( \psi(M_i + \lvert G \rvert \alpha)
                     - \psi(\lvert G \rvert \alpha) \bigr)} ,
    \end{equation}
    %
    interleaved with the sampler, which carries an assignment state per sample and sweeps them in
    turn. The update is derived by bounding the Dirichlet log-likelihood and converges monotonically,
    so it needs no step size either --- the same property that recommends the iteration of
    Section \ref{inference} over an ascent method. Alternatively $\alpha$ is given a hyperprior and
    sampled with the rest, which costs less thought and about as much time.

    \todo{What $m$ has to be for the estimate to mean anything, and what to do below it. A single
    sample is stated above as uninformative; a handful is not obviously better, and the paper should
    say where the boundary lies rather than leave a reader to assume any cohort will do. The
    fallback when $m$ is too small is a fixed $\alpha$ chosen and declared, which is honest, rather
    than an estimated one that is noise wearing a number.}

    \subsection{Counts that are not integers}
    A sampler over occurrences needs occurrences, and the
    corrected $\hat{c}_n$ of Section \ref{error} are not whole numbers. Rounding them is the obvious
    repair and the wrong one: it distorts the correction most where counts are small, which is where
    a rare taxon lives. It is not needed. The correction is a single factor $(1 - e_b)^k$ applied to
    every $c_n$ alike, and what the prior weighs itself against is the ratio of counts to
    concentration, not their absolute size. Sampling the raw integer occurrences $c_n$ under
    %
    \begin{equation}\label{smoothyalphaeff}
        \alpha_{\mathrm{eff}} \;=\; \alpha \, (1 - e_b)^k
    \end{equation}
    %
    leaves that ratio exactly where $\hat{c}_n$ with $\alpha$ would have put it, and keeps the state
    integral.

    That holds only while the correction is a single factor, and Section \ref{error} has since given
    it a second one that is not. The consistency weight $w(r_n)$ differs from node to node, so no
    rescaling of $\alpha$ can reproduce it: the weights change the counts relative to \emph{each
    other} and not merely their common size. A sampler would have to carry them as per-occurrence
    weights, which is no longer collapsed Gibbs sampling but a weighted variant of it, or fall back on
    rounding $w(r_n) c_n$ with the loss that was the objection in the first place. The iteration of
    Section \ref{inference} has no such difficulty --- a weighted likelihood is what it already
    maximises --- which is one more reason the prior sits in this appendix and not in the method.

    \todo{What the sampler costs and how it is run: burn-in, how many sweeps are retained, and what
    is used to decide that it has mixed. Report it rather than choosing it silently --- a sampler with
    an undeclared burn-in is not a reproducible measurement. The degenerate taxa of
    Section \ref{pgn} are where mixing will be worst, since it is precisely there that the posterior
    has a ridge rather than a peak, so they are the place to look.}

    \end{appendices}

    \backmatter

    \bmhead{Acknowledgements}

    \section*{Declarations}

    \bmhead{Availability of data and materials}
    \todo{}

    \bmhead{Competing interests}
    The authors declare that they have no competing interests.

    \bmhead{Funding}
    \todo{}

    \bmhead{Authors' contributions}
    \todo{}

    \bibliography{sn-bibliography}

\end{document}
